O estoque ambiental afetado pelas decisões de uso da terra pode modificar as condições biofísicas que sustentam a produção e, por essa via, alterar a resposta da agricultura a perturbações climáticas. Este artigo propõe investigar se o estoque de vegetação nativa previamente existente no território modifica, em magnitude economicamente relevante, a resposta da produtividade agrícola a choques de precipitação, e se essa relação apresenta heterogeneidade entre Amazônia, Cerrado e Pantanal em Mato Grosso. Para tanto, será construído um painel territorial, com período candidato de 2001 a 2024, combinando rendimento de soja da Produção Agrícola Municipal, cobertura da terra do MapBiomas e precipitação do CHIRPS. A estratégia econométrica considerará a interação entre choques climáticos e vegetação anterior à janela agrícola, efeitos fixos e diferenças nas condições produtivas locais. Não se impõe, de antemão, que a vegetação necessariamente amorteça o choque. A interpretação dos resultados dependerá da harmonização territorial, do suporte amostral e das condições de identificação, particularmente da possível seleção dos estoques vegetais e da influência da cobertura sobre o clima local. Nesta etapa de qualificação, apresentam-se a fundamentação e o desenho empírico; não há resultados econométricos estimados.

\medskip
\noindent\textbf{Palavras-chave:} capital natural; produtividade agrícola; choques climáticos; vegetação nativa; Mato Grosso.
